% -----------------------------------------------------------------------------------------------------
\section{算符乘积展开的严格结果}\label{sec:front-ope}

OPE 是最古老的因子化定理，也是唯一有真正数学物理意义上严格结果的一个。Zimmermann 在 BPHZ 重整化的标量理论中证明了复合算符乘积的短距离展开，在插入 Green 函数的意义下成立，保留足够多算符后余项比任意给定幂次更快地趋于零\cite{Zimmermann1973}；规范理论的规范不变版本另需 BRST 与算符混合的分类（\cref{sec:front-brst}）。2006 年以后的结果在特定理论中把它做成自洽的严格证明，并增加了新的内容（弯曲时空、收敛性、结合律、递推关系）。问题是：在哪个理论中、在什么意义下（渐近还是收敛）证明了什么。

\begin{proposition}[标量场论中 OPE 的收敛、结合律与弯曲时空 \tagP]\label{prop:front-ope-scalar}
在欧氏 $g\varphi^4_4$（有质量与无质量）中，在每个固定圈数下：(a) OPE 在旁观者场具有紧支动量的态中对每个有限 $|x|$ 收敛——有质量情形即\cref{thm:power-ope-convergence}（界见\cite[式~(7)]{Hollands:2011gf}），无质量情形用处理例外动量奇异性的新技术（“树优势”）得到同样结论\cite{Holland:2014pna}；(b) OPE 系数满足只含 OPE 系数本身的紫外有限的递推公式（对耦合常数求导），可由自由场系数递推构造\cite{Holland:2014ifa}，并在所有阶满足强结合律，两点系数决定 $n$ 点系数（“相干定理”）\cite{Holland:2015tia}；(c) 在任意整体双曲的 Lorentz 弯曲时空上，微扰相互作用标量场论有任意阶的渐近 OPE，余项在任意 Hadamard 态中于短距离趋于零，系数局域协变、满足结合律、重整化群方程与微局域条件\cite{Hollands:2006ag}。（(a) 的无质量情形与 (b)(c) 均据各文摘要。本命题是标量模型中的陈述。）
% VF-FRO-4
\uses{thm:power-ope-convergence}
\end{proposition}

\begin{proofsketch}
(a) 的架构（\cref{thm:power-ope-convergence}）：带红外截断 $\Lambda$ 与紫外截断 $\Lambda_0$ 的 Wilson--Wegner--Polchinski 流方程；对 OPE 余项按圈数与外腿数作归纳估计，估计对 $\Lambda_0$ 一致；所用的加细归纳假设能自我再现；附带结果是固定圈数下有效作用量梯度展开的收敛性。完整论证见\cite{Hollands:2011gf,Holland:2014pna}。
\end{proofsketch}

\paragraph{地位与缺口。}[P]：作为与 QCD 相关的陈述，它只在标量模型中成立（结果发表于经审稿的数学物理期刊）。原文的保留：收敛性依赖于态（旁观者动量紧支；态的高能行为越差收敛越差），圈数级数本身不在讨论之内；在 Minkowski 空间中只可期望类空间隔下收敛，因此对 DIS 的光锥展开，除矩之外得不出任何结论。

\begin{proposition}[纯 Yang--Mills 理论中 OPE 在规范不变算符上闭合 \tagP]\label{prop:front-ope-ym}
在纯 Yang--Mills 理论（半单规范群）中：(a) 微扰 Yang--Mills 理论在任意整体双曲时空上作为可观测量的形式幂级数代数存在（辅助规范固定代数的 BRST 上同调；重整化的 Ward 恒等式保证相互作用 BRST 流守恒与 BRST 荷幂零），OPE 在规范不变算符上闭合，重整化群流不离开它们\cite{Hollands:2007zg}（据摘要与适用范围；物质场只被讨论，手征费米子可能有反常，含夸克的 QCD 不在其中）；(b) 对任意无质量、超表观可重整的理论，在任意固定阶，可选取重整化条件使 OPE 存在，且只保留维数低于阈值 $\Delta$（原文记为 $D$）的算符时余项满足
\begin{equation}\label{eq:front-fh}
\lim_{\tau\to0}\tau^{[\mathcal O_A]-\Delta+\delta}\Bigl[\bigl\langle\mathcal O_{A_1}(\tau x_1)\cdots\mathcal O_{A_s}(\tau x_s)\bigr\rangle-\sum_{[\mathcal O_B]<\Delta}\mathcal C^B_{A_1\cdots A_s}(\tau x)\bigl\langle\mathcal O_B(\tau x_s)\bigr\rangle\Bigr]=0\qquad(\forall\delta>0),
\end{equation}
其中 $[\mathcal O_A]=\sum_k[\mathcal O_{A_k}]$ 是各算符的工程维数之和；该展开在由无例外动量的涂抹场生成的态中亦然；对无质量可重整规范理论，OPE 系数满足 BRST Ward 恒等式，前提是 $E(4)$ 等变上同调类 $H^{1,4}_{E(4)}(\hat s|d)$ 为零——纯 Yang--Mills（半单群）满足，“but not, for instance, in the presence of chiral fermions in non-real representations”；系数满足递推关系与结合律\cite[第~1.1~节]{Frob:2016mzv}（预印本）。（本命题是纯 Yang--Mills 理论中的微扰陈述。）
\uses{def:prelim-qcd, def:prelim-brst}
\end{proposition}

\paragraph{地位与缺口。}[P]：作为 QCD 陈述（不含夸克）。(b) 的证明仅见于预印本，它所用的界（\cite{Frob:2015uqy}）只对纯 Yang--Mills 建立。对含夸克的 QCD，微扰 OPE 在本书中仍是全阶论证（\cref{thm:dis-ope}，[A]）。

\begin{proposition}[流方程（Wilson--Polchinski）重整化证明 \tagP]\label{prop:front-flow-renorm}
对下列理论，用流方程证明了带显式界的微扰可重整性（均据各文摘要）：有质量 $\varphi^4_4$\cite{Keller:1990ej}；含电子的 QED\cite{Keller:1991cj,Keller:1995qn}；自发破缺的 SU(2) Yang--Mills--Higgs，并恢复 Slavnov--Taylor 恒等式\cite{Kopper:1999qx,Kopper:2009zw}；无质量 $\varphi^4_4$ 的全阶一致动量界\cite{Guida:2011np}（预印本）；采用物理重整化条件、控制无质量关联函数红外行为的纯 SU(2) Yang--Mills\cite{Efremov:2017sqi}；半单群欧氏 Yang--Mills 含任意复合算符的情形，包括 BRST Ward 恒等式、树加权的界，紫外与红外均受控\cite{Frob:2015uqy}。（本命题是所列模型中的微扰陈述。）
\uses{def:prelim-qcd, def:prelim-brst}
\end{proposition}

\paragraph{地位与缺口。}[P]：作为 QCD 陈述。据我们截至 2026 年 10 月的文献检索（INSPIRE：“flow equation” 与 QCD／费米子），尚未见含夸克的 QCD 的自洽流方程证明（检索并不穷尽）。

\begin{conjecture}[Yang--Mills 与 QCD 中 OPE 的收敛 \tagC]\label{conj:front-ope-convergence}
对纯 Yang--Mills 与含夸克的 QCD，在每一固定圈数下，规范不变复合算符的欧氏 OPE 在动量紧支的态中对每个有限 $|x|$ 收敛（\cref{thm:power-ope-convergence}的类比）。
\uses{thm:dis-ope, prop:front-ope-ym}
\end{conjecture}

\paragraph{地位与缺口。}[C]：Fröb 与 Holland 认为原则上没有障碍，但需要更精细的界\cite{Frob:2016mzv}（预印本）；即便成立，它也是欧氏（类空）的微扰陈述。

\begin{remark}[对含夸克的 QCD，本书能断言什么]\label{rem:front-ope-qcd-status}
(1) 含规范不变算符的全阶渐近 OPE：本书记为全阶论证（\cref{thm:dis-ope}，[A]）。把 Zimmermann 正规乘积\cite{Zimmermann1973}、BRST、Joglekar--Lee 分类\cite{Joglekar:1975nu}与代数重整化\cite{Piguet:1995er}（教科书形式见\cite{Collins1984}）组合起来即得 [R] 的说法，本书不采纳：组合本身没有以单一自洽证明的形式写出。(2) 收敛性只对 $\varphi^4_4$ 证明。(3) DIS 的 Minkowski 光锥 OPE 来自欧氏 OPE 加上矩的解析性与色散关系（第七章）。(4) 以上都不是非微扰结论；作为量子场论公理的 OPE（Hollands--Wald 纲领）只是受这些结果启发，并非由它们推出（\cref{asm:intro-np-ope}）。不加理论与阶数限定地说“OPE 已被严格证明”是过度表述。
\end{remark}
